% Generated draft: requires mathematical and source review.
\documentclass[11pt]{article}
\usepackage[margin=25mm]{geometry}
\usepackage{fontspec}
\setmainfont{Latin Modern Roman}
\usepackage{amsmath,amssymb,mathtools,unicode-math}
\setmathfont{Latin Modern Math}
\usepackage{enumitem,needspace,xurl,hyperref}
\setlength{\parindent}{0pt}
\setlength{\parskip}{5pt}
\setlength{\emergencystretch}{3em}
\newcommand{\N}{\mathbb{N}}
\newcommand{\Z}{\mathbb{Z}}
\newcommand{\Q}{\mathbb{Q}}
\newcommand{\R}{\mathbb{R}}
\newcommand{\C}{\mathbb{C}}
\newcommand{\F}{\mathbb{F}}
\newcommand{\A}{\mathbb{A}}
\newcommand{\PP}{\mathbb{P}}
\newcommand{\E}{\mathbb{E}}
\newcommand{\eps}{\varepsilon}
\newcommand{\dd}{\,\mathrm{d}}
\begin{document}
\section*{UnsolvedMath ID: 30002843}
\textbf{Source problem code:} catalogue entry
\subsection{Definitions and mathematical statement}
For a vector \(v\in\C^d\), let \(v^*\) denote its conjugate transpose, so that \(vv^*\) is the rank-one positive semidefinite matrix with entries \(v_j\overline{v_k}\). Write \(I_d\) for the identity matrix, use the Euclidean norm on vectors, and use the induced operator norm on matrices. A random vector has finite support if it takes values in a finite set with probability one.

The supplied statement asserts the following implication. Let \(\varepsilon>0\), and let \(v_1,\ldots,v_m\) be independent random vectors with finite support in \(\C^d\), defined on a common probability space. Suppose that
\[
\sum_{i=1}^{m}\E(v_i v_i^*)=I_d,
\qquad
\E\|v_i\|^2\leq\varepsilon
\quad\text{for every }i\in\{1,\ldots,m\}.
\]
Then
\[
\Pr\!\left(
\left\|\sum_{i=1}^{m}v_i v_i^*\right\|
\leq(1+\sqrt{\varepsilon})^2
\right)>0.
\]
The record uses \(m\) and \(d\) as finite family-size and dimension parameters without explicitly specifying whether zero values are permitted; no additional boundary convention is asserted here.

\textit{Unresolved source qualifications.} The supplied verification field identifies this passage as a proved Marcus--Spielman--Srivastava theorem, rather than an open problem. However, the registry labels the record \texttt{open} and its statement status \texttt{unrecoverable}. Its dated literature assessment says that the conclusion is missing, although both supplied statement fields contain the conclusion displayed above. These conflicting annotations are preserved as source qualifications; no separate open question is supplied, and no independent mathematical or literature verification is asserted.
\subsection{Short English statement}
The supplied implication says that independent complex random vectors with finite support, whose expected outer products sum to the identity and whose individual expected squared lengths are at most \(\varepsilon\), have positive probability of producing a sum of outer products with operator norm at most \((1+\sqrt{\varepsilon})^2\). The record identifies this as a proved Marcus--Spielman--Srivastava theorem, but its conflicting open-problem and extraction annotations require review.
\subsection{Sources}
\begin{itemize}
\item Local UnsolvedMath registry: \texttt{data/unsolvedmath/problems.json}, permanent ID \(30002843\), exact source problem code \texttt{catalogue entry}, title ``Marcus--Spielman--Srivastava Theorem Extraction.'' Collection: Oberwolfach Reports Open Problems. No individual proposer is supplied.
\item Supplied report citation: \emph{Arbeitsgemeinschaft: The Kadison-Singer Conjecture} (2016), Oberwolfach Reports, pp.~927--969, DOI \texttt{10.4171/owr/2015/17}. Supplied source URL and literature reference labelled ``original OWR report'': \url{https://doi.org/10.4171/owr/2015/17}. The record separately gives proposed year 2015; the relationship between these year labels is not resolved here.
\item Attribution: the supplied statement-verification field calls the passage a proved Marcus--Spielman--Srivastava theorem. No separate original theorem citation, theorem number, or page locating this passage is supplied.
\item Literature assessment dated \texttt{2026-08-21}: supplied dated curation, not independently verified. It labels the record open under a queue rule because it says the theorem conclusion is absent. That explanation conflicts with the complete implication in the supplied statement fields. The background expressly requests independent verification of time-sensitive status and cited claims.
\end{itemize}
\end{document}
